% related_work.tex -- Related Work for the IRAI 2027 paper (IEEEtran, conference mode).
% All \cite keys exist in literature/references.bib. 35 distinct references are cited here
% (the table cites 8 of them). Prose is ~565 words excluding citations.
% Space note: IRAI counts references inside the 6-page limit. If space is tight, the first
% candidates to drop are boisvert2024workarenapp, shi2025progent, liu2023verifiability,
% breck2017mltestscore, cui2025orbench and gupta2026reliabilitybench (each paragraph's
% argument survives without them).
% The comparison table is in related_work_table.tex (needs \usepackage{booktabs,amssymb}).

\section{Related Work}
\label{sec:related}

\textbf{Enterprise agent benchmarks.}
Agent benchmarks have moved from single-platform knowledge work~\cite{boisvert2024workarenapp} to simulated companies~\cite{xu2025theagentcompany} and policy-constrained tool use, where repeated trials expose large run-to-run inconsistency~\cite{yao2025taubench}. EnterpriseBench~\cite{vishwakarma2025enterprisebench} adds departmental personas, fragmented sources and access-control hierarchies, and CRMArena-Pro~\cite{huang2026crmarenapro} reports near-zero confidentiality awareness that prompting improves, often at a cost to task success. Closer to our source conditions, WorkSurface-Bench~\cite{liang2026worksurfacebench} scores knowledge-surface routing separately from answer quality, EnterpriseRAG-Bench~\cite{sun2026enterpriseragbench} injects conflicting and absent information, and LayerRAG-Bench~\cite{shams2026layerrag} injects stale-index and permission-denied faults and shows that groundedness-only scoring produces false positives. ST-WebAgentBench~\cite{levy2026stwebagentbench} similarly finds that completion under policy falls well below nominal completion. These benchmarks evaluate one deployed configuration and mostly score task completion. We instead hold the agent fixed, vary release-stage configuration and enforcement mode, and add a relevant-but-prohibited source condition that none of them models.

\textbf{Access control and policy enforcement.}
Access rules stated in the prompt are followed unreliably. OrgAccess~\cite{sanyal2026orgaccess} shows that LLMs fail role-based permissions as rules compose, often answering over-cautiously; ACCESS DENIED INC~\cite{fazlija2025accessdenied} tests whether models refuse unauthorized employee requests while still answering legitimate ones; and the system/user prompt split does not yield a dependable instruction hierarchy~\cite{geng2026controlillusion}. Enterprise privacy benchmarks report frequent leakage and a utility--privacy trade-off~\cite{fu2026ciwork,baek2026leakage}. System-level alternatives enforce policy outside the model through capabilities~\cite{debenedetti2025camel}, tool-call privilege policies~\cite{shi2025progent} or reference monitors~\cite{palumbo2026policy}. Direct comparisons favour architectural enforcement over prompt allowlists for tool access~\cite{uppala2026promptsdontprotect}, and retrieval-time attribute-based gating removes cross-tenant leakage~\cite{arceo2026securing}. We therefore treat the advantage of tool-level enforcement as an expected effect to quantify, not a discovery. We measure it for record retrieval and disclosure, jointly with accuracy, evidence support and over-restriction, and test whether it holds across release-stage packages and under permission drift.

\textbf{Attribution, staleness and abstention.}
Attribution research defines support against an identified source~\cite{rashkin2023ais} and automates claim-level citation checking~\cite{gao2023alce}. Audits find that fluent answers are often only partly supported by their citations~\cite{liu2023verifiability}, and Trust-Score~\cite{song2025trustscore} shows that parametric knowledge inflates correctness relative to grounded citation and refusal. A correct citation is also not necessarily the evidence the model relied on~\cite{wallat2025correctness}, which motivates our trace-based check that the cited source was actually retrieved. GaRAGe~\cite{sorodoc2025garage} adds deflection and time-sensitive questions, HoH~\cite{ouyang2025hoh} shows that outdated documents mislead RAG even when current ones are retrieved, and RGB~\cite{chen2024rgb} tests rejection when no retrieved document holds the answer, while over-refusal benchmarks show that stricter behaviour also rejects benign requests~\cite{cui2025orbench}. In all of these, any retrieved document is admissible evidence, so a claim fully supported by a stale or prohibited record counts as supported. We score Source Appropriateness (the authoritative permitted source was retrieved and used) separately from Supported Answer Rate and Access-Boundary Compliance, and we separate correct refusal on no-permitted-source items from over-restriction.

\textbf{Lifecycle evaluation.}
Production ML practice gates promotion on tests and post-release monitoring~\cite{breck2017mltestscore}, and EDDOps~\cite{xia2024eddops} proposes evaluation that spans offline development and online operation of LLM agents. Empirical lifecycle evidence remains narrow: model-version drift on static tasks~\cite{chen2024chatgptdrift}, a longitudinal release-gate case study of one internal system~\cite{maiorano2026qualitygate}, stress tests with injected tool faults such as schema drift~\cite{gupta2026reliabilitybench}, and practitioner reports that reliability is the main production bottleneck~\cite{pan2026measuring}. The NIST generative-AI profile warns that pre-deployment testing may not match deployment contexts~\cite{nist2024genai600} but does not quantify the gap. We contribute a controlled factorial comparison of stage-modelled configuration packages with documented per-package differences, injected document and permission drift, and mixed-effects estimates with scenario random effects, so that each stage contrast is reported with its uncertainty rather than as a single readiness score.
